% Compile with: uplatex file.tex && dvipdfmx file.dvi
%           or: latexmk -uplatex file.tex
\documentclass[dvipdfmx]{article}
\usepackage{amsmath, amssymb}
\usepackage{tikz}

\title{scenario\_utility\_maximization}
\date{\today}

\begin{document}
\maketitle

Maximise the Cobb–Douglas utility $U(x_1,x_2)=x_1^ax_2^b$ subject to $p_1x_1+p_2x_2=m$. First declare the subscripted symbols as NAMES — undeclared, they resolve at operation time as the component index $x(1)$ and the first-order conditions come out wrong.

declare $x_1$ as a scalar

\[
  x_1
\]

declare $x_2$ as a scalar

\[
  x_2
\]

declare $p_1$ as a scalar

\[
  p_1
\]

declare $p_2$ as a scalar

\[
  p_2
\]

Register the utility function.

$U = x_1^a x_2^b$

\begin{equation}
  U = x_1^a x_2^b
  \label{eq:E1}
\end{equation}

The researcher writes the Lagrangian, and names the multiplier (the system need not auto-number one).

$L = U - \lambda (p_1 x_1 + p_2 x_2 - m)$

\begin{equation}
  L = U - \lambda (p_1 x_1 + p_2 x_2 - m)
  \label{eq:E2}
\end{equation}

For the first-order conditions, differentiate the Lagrangian with respect to each variable.

take the partial derivative of E2 with respect to $x_1$

\begin{equation}
  \frac{\partial L}{\partial x_1} = x_{1}^{a-1} x_{2}^{b} a-p_{1} \lambda
  \label{eq:E3}
\end{equation}

take the partial derivative of E2 with respect to $x_2$

\begin{equation}
  \frac{\partial L}{\partial x_2} = x_{1}^{a} x_{2}^{b-1} b-p_{2} \lambda
  \label{eq:E4}
\end{equation}

take the partial derivative of E2 with respect to $\lambda$

\begin{equation}
  \frac{\partial L}{\partial \lambda} = m-p_{2} x_{2}-p_{1} x_{1}
  \label{eq:E5}
\end{equation}

Set each right-hand side to zero to get the first-order conditions.

set the right-hand side of E3 to 0

\begin{equation}
  x_{1}^{a-1} x_{2}^{b} a-p_{1} \lambda = 0
  \label{eq:E6}
\end{equation}

set the right-hand side of E4 to 0

\begin{equation}
  x_{1}^{a} x_{2}^{b-1} b-p_{2} \lambda = 0
  \label{eq:E7}
\end{equation}

set the right-hand side of E5 to 0

\begin{equation}
  m-p_{2} x_{2}-p_{1} x_{1} = 0
  \label{eq:E8}
\end{equation}

Do the comparative statics FIRST. Write the demand relation explicitly as $F(x_1,m)=am-(a+b)p_1x_1=0$. This step must precede the solve: solving binds $x_1$ on the CAS side, which collapses $F$ identically to 0.

$F = a m - (a+b) p_1 x_1$

\begin{equation}
  F = a m - (a+b) p_1 x_1
  \label{eq:E9}
\end{equation}

Obtain the two partial derivatives the implicit function theorem needs.

take the partial derivative of E9 with respect to m

\begin{equation}
  \frac{\partial F}{\partial m} = a
  \label{eq:E10}
\end{equation}

take the partial derivative of E9 with respect to $x_1$

\begin{equation}
  \frac{\partial F}{\partial x_1} = -p_{1} \left(b+a\right)
  \label{eq:E11}
\end{equation}

The researcher writes the implicit function theorem $dx_1/dm=-F_m/F_{x_1}$, pulling in the previous results by expression reference.

$\frac{dx_1}{dm} = -\frac{E10}{E11}$

\begin{equation}
  \frac{d}{d m} x_{1}=\frac{a}{p_{1} \left(b+a\right)}
  \label{eq:E12}
\end{equation}

Solve the first-order conditions for the demand functions. (Written with `and` rather than a comma list — see open\_items SOLVE-LIST-COMMA-EN.)

solve E6, E7 and E8 for $x_1$, $x_2$, $\lambda$

\begin{equation}
  x_{1} = \frac{a m}{p_{1} b+p_{1} a}
  \label{eq:E13}
\end{equation}

\begin{equation}
  x_{2} = \frac{b m}{p_{2} b+p_{2} a}
  \label{eq:E14}
\end{equation}

\begin{equation}
  \lambda = {{b \left(\frac{a m}{p_{1} b+p_{1} a}\right)^{a} \left(\frac{b m}{p_{2} b+p_{2} a}\right)^{b}+a \left(\frac{a m}{p_{1} b+p_{1} a}\right)^{a} \left(\frac{b m}{p_{2} b+p_{2} a}\right)^{b}}\over{m}}
  \label{eq:E15}
\end{equation}

Differentiate the solved demand function with respect to $m$ directly, and confirm it agrees with the implicit-function-theorem result. The macro returned one answer with nothing to compare it against.

take the partial derivative of E13 with respect to m

\begin{equation}
  \frac{\partial x_1}{\partial m} = \frac{a}{p_{1} b+p_{1} a}
  \label{eq:E16}
\end{equation}

\end{document}